\section{Summary}

% =========================================================
% (A) Organizing frame: training & execution paradigms
% =========================================================
\begin{frame}
\frametitle{Chapter 9 --- The Big Picture}
\framesubtitle{Training \& execution paradigms}

Across all of Chapter 9, algorithms differ in \emph{what information is shared,
and when.} \textbf{Centralized information} is anything beyond an agent's own
observation --- parameters, gradients, observations, or actions of others.

\vspace{0.6em}
\centering
\begin{tikzpicture}[
    every node/.style={font=\small},
    para/.style={
        rectangle, rounded corners=4pt, very thick,
        align=center, inner sep=6pt,
        text width=3.0cm, minimum height=1.9cm
    },
    cte/.style={para, draw=blue!55!black, fill=blue!8},
    dte/.style={para, draw=green!55!black, fill=green!8},
    ctde/.style={para, draw=orange!85!black, fill=orange!10}
]

\node[dte] (dte) at (0,0)
    {\textbf{Decentralized training \& execution} \\ {\scriptsize agents fully isolated; learn from local observations only}};

\node[ctde] (ctde) at (4.5,0)
    {\textbf{Centralized training, decentralized exec.\ (CTDE)} \\ {\scriptsize share info at \emph{training}; act on local obs.\ at execution}};

\node[cte] (cte) at (9.0,0)
    {\textbf{Centralized training \& execution} \\ {\scriptsize share info throughout both training \emph{and} execution}};

% spectrum arrow underneath
\draw[{Latex}-{Latex}, thick, black!55]
    ($(dte.south)+(0,-0.35)$) -- ($(cte.south)+(0,-0.35)$);
\node[font=\scriptsize, text=black!60] at ($(dte.south)+(0,-0.65)$) {less sharing};
\node[font=\scriptsize, text=black!60] at ($(cte.south)+(0,-0.65)$) {more sharing};

\end{tikzpicture}

\vspace{0.5em}
{\small CTDE is the popular middle ground: exploit centralized information to
\emph{learn} better, keep policies deployable with only local information.}

\end{frame}

\begin{frame}
\frametitle{Chapter 9 --- Four Recurring Challenges}

Deep MARL is largely the story of \emph{four} core challenges --- each section
attacks one or more of them:

\vspace{0.4em}
\begin{columns}[T,onlytextwidth]

\begin{column}{0.48\textwidth}
\begin{block}{Non-stationarity}
\small Others learn and change $\Rightarrow$ a moving target.
\emph{Centralized critics (9.4), agent modeling (9.6).}
\end{block}

\begin{block}{Multi-agent credit assignment}
\small Who caused the shared reward?
\emph{Value decomposition --- VDN, QMIX (9.5).}
\end{block}
\end{column}

\begin{column}{0.48\textwidth}
\begin{block}{Scaling \& sample efficiency}
\small Cost grows with the number of agents.
\emph{Parameter \& experience sharing (9.7).}
\end{block}

\begin{block}{Equilibrium selection}
\small Which joint policy to converge to?
\emph{Self-play (9.8), population-based training (9.9).}
\end{block}
\end{column}

\end{columns}

\vspace{0.8em}
\centering
\emph{Partial observability} cuts across all four --- handled throughout via
observation histories $h_i^t$ and centralized information $z^t$.



\end{frame}

% =========================================================
% (B) The section-by-section table (refined)
% =========================================================
\begin{frame}
\frametitle{What did we discuss in Chapter 9?}

\small

\centering
\begin{tabularx}{\textwidth}{
>{\bfseries}l
>{\raggedright\arraybackslash}l
>{\raggedright\arraybackslash}X
}
\toprule
Section & Technique & Challenge addressed \\
\midrule

9.3 Independent Learning
&
Treat others as environment
&
Baseline --- exposes non-stationarity
\\

9.4 Policy Gradients
&
Centralized critics
&
Credit assignment, non-stationarity (training)
\\

9.5 Value Decomposition
&
VDN, QMIX
&
Multi-agent credit assignment, at scale
\\

9.6 Agent Modeling
&
Deep JAL-AM, representations
&
Reason about \& adapt to other agents; generalize to unseen states
\\

9.7 Homogeneous Agents
&
Parameter \& experience sharing
&
Scaling, sample efficiency
\\

9.8 Self-Play
&
MCTS, AlphaZero
&
Equilibrium-finding, 2-player zero-sum
\\

9.9 Population-Based Training
&
PSRO, AlphaStar
&
Equilibrium-finding \& robustness, general competitive
\\

\bottomrule
\end{tabularx}

\vspace{0.6em}
{\footnotesize All build on the same paradigms (decentralized / CTDE /
centralized) --- they differ in \emph{what} is shared and \emph{which}
challenge it targets.}

\end{frame}

% =========================================================
% (C) Closing big-picture: challenges -> techniques
% =========================================================
